\section[Hex - Dez]{Hex $\to$ Dez}
\begin{ctabular}{l|cccccc}
    \textbf{Hex} & A & B & C & D & E & F \\
    \textbf{Dec} & 10 & 11 & 12 & 13 & 14 & 15 \\
    \textbf{Bin} & 1010 & 1011 & 1100 & 1101 & 1110 & 1111
\end{ctabular}\vspace{-1em}

\section{Instruction Set Architecture (ISA)}
\label{subsec:isa_differences}

Die Weiterentwicklung der ARM-Architekturen hat zur Koexistenz verschiedener Befehlssätze geführt.

\begin{itemize}
    \item \textbf{ARM (32-bit) (Veraltet):} 32 Bit, maximale Leistung, aber sie belegen mehr Speicherplatz.
    \item \textbf{Thumb (16-bit):} Untermenge 16 Bit,
    Ideal, um den Speicher-Footprint zu reduzieren (ca. 30\% Einsparung)
    \item \textbf{Thumb-2 (16/32-bit):} variable Länge (16 und 32 Bit). Ermöglicht eine Code-Dichte ähnlich zu Thumb 
    mit vergleichbarer Leistung wie der 32-Bit-ARM-Befehlssatz.
\end{itemize}

\section{Logik der ALU und Program Status Register}
\label{sec:alu_logic}

Die Arithmetisch-Logische Einheit (ALU) arbeitet ausschließlich auf binären Bitfolgen, 
ohne die Semantik des Vorzeichens zu berücksichtigen (Interpretation \textit{signed} vs \textit{unsigned}).

\subsection{Flag-Mechanismus (APSR)}
\label{subsec:apsr_flags}

\begin{itemize}
    \item \textbf{N (Negative):} Spiegelt das höchstwertige Bit (MSB) des Ergebnisses wider. Wenn $MSB=1$, wird das Flag auf 1 gesetzt.
    \item \textbf{Z (Zero):} Auf 1 gesetzt, wenn alle Bits des Ergebnisses gleich 0 sind.
    \item \textbf{C (Carry):} Zeigt einen Übertrag über das höchstwertige Bit hinaus an. 
\begin{itemize}
    \item Entscheidend für Operationen \textit{unsigned}. \crd{Für die Subtraktion invertiert!}
    \item \textbf{Subtraktions-Mechanismus:} $(A + \text{NOT}(B) + 1)$. \quad Zweierkomplement.
\begin{itemize}
        \item $C = 1 \implies A \ge B$ (Kein Borrow / Not-Borrow).
        \item $C = 0 \implies A < B$ (Borrow erfolgt / Borrow).
    \end{itemize}
\end{itemize}
    \item \textbf{V (oVerflow):} Zeigt einen Vorzeichenfehler in der Rechnung im Zweierkomplement an 
    und ist entscheidend für Operationen \textit{signed}.
\begin{itemize}
        \item $(+) + (+) = (-)$ Unmöglich
        \item $(-) + (-) = (+)$ Unmöglich
    \end{itemize}
\end{itemize}

\subsubsection{Interpretation des Vorzeichens und bedingte Sprünge}
\label{subsec:sign_interpretation}

Da die ALU gleichzeitig sowohl das Flag $C$ als auch das Flag $V$ aktualisiert,
liegt die Unterscheidung zwischen der Variable \texttt{int} und \texttt{uint} nicht in der Hardware, 
sondern in der Software.

